\documentclass[11pt,a4paper]{article}

% ── Packages ──────────────────────────────────────────────────────────────
\usepackage[utf8]{inputenc}
\usepackage[T1]{fontenc}
\usepackage{amsmath,amssymb,amsthm}
\usepackage{mathtools}
\usepackage{bm}
\usepackage{geometry}
\geometry{margin=1in}
\usepackage{hyperref}
\hypersetup{colorlinks=true,linkcolor=blue,citecolor=blue,urlcolor=blue}
\usepackage{listings}
\usepackage{xcolor}
\usepackage{titlesec}
\usepackage{enumitem}
\usepackage{booktabs}
\usepackage{fancyhdr}
\usepackage{graphicx}
\usepackage{longtable}
\usepackage{caption}
\usepackage{subcaption}

% ── Code listing style ───────────────────────────────────────────────────
\definecolor{codebg}{RGB}{248,248,248}
\definecolor{codeframe}{RGB}{200,200,200}
\definecolor{codecomment}{RGB}{0,128,0}
\definecolor{codestring}{RGB}{163,21,21}
\definecolor{codekeyword}{RGB}{0,0,200}
\definecolor{diffadd}{RGB}{0,128,0}
\definecolor{diffrem}{RGB}{200,0,0}

\lstdefinestyle{pythonstyle}{
    language=Python,
    backgroundcolor=\color{codebg},
    frame=single,
    rulecolor=\color{codeframe},
    basicstyle=\ttfamily\small,
    keywordstyle=\color{codekeyword}\bfseries,
    commentstyle=\color{codecomment}\itshape,
    stringstyle=\color{codestring},
    showstringspaces=false,
    breaklines=true,
    breakatwhitespace=false,
    tabsize=4,
    captionpos=b,
    numbers=left,
    numberstyle=\tiny\color{gray},
    numbersep=8pt,
    xleftmargin=1.5em,
    framexleftmargin=1.5em,
    aboveskip=1em,
    belowskip=1em,
    morekeywords={self,True,False,None,as,with,lambda},
    literate={→}{$\rightarrow$}1 {←}{$\leftarrow$}1,
}
\lstset{style=pythonstyle}

% ── Header / Footer ──────────────────────────────────────────────────────
\pagestyle{fancy}
\fancyhf{}
\rhead{DRP Modifications Report}
\lhead{Optimizing the Optimizer --- Cahn--Hilliard Extension}
\cfoot{\thepage}

% ── Title ─────────────────────────────────────────────────────────────────
\title{%
  \textbf{DRP Modifications Report}\\[6pt]
  \large Spectral Reference Solution, SSBroyden1 Bug Fix,\\
  Negative Curvature Guard, and RAdam Warm-Up\\[6pt]
  \normalsize Extending \emph{Optimizing the Optimizer for Data-Driven PINNs and KANs}\\
  (Kiyani et al., CMAME \textbf{446}:118308, 2025)
}
\author{Maria Emilia Ciasullo}
\date{February 26, 2026}

% ══════════════════════════════════════════════════════════════════════════
\begin{document}
\maketitle
\tableofcontents
\newpage

% ======================================================================
\section{Introduction}
\label{sec:intro}
% ======================================================================

This document records all modifications made to the codebase of
Kiyani et al.\ \cite{kiyani2025} during the Directed Research Project
(DRP) on extending their quasi-Newton optimizer family to the
\textbf{Cahn--Hilliard equation} (Model~B dynamics).  Every change is
fully commented in the source files with \texttt{DRP MODIFICATION
START/END} markers so that the original paper code can be recovered at
any time.

The work described here covers three priorities:

\begin{enumerate}[label=\textbf{P\arabic*}]
  \item \textbf{Spectral Reference Solution} --- build a ground-truth
        solver and compute relative $\ell_2$ errors for all 17
        experiments (\S\ref{sec:spectral}).
  \item \textbf{SSBroyden1 Bug Fix \& Line Search Audit} --- diagnose
        and repair numerical failures in the SSBroyden1 Hessian update,
        add a negative-curvature guard affecting all variants
        (\S\ref{sec:ssbroyden1}).
  \item \textbf{RAdam Warm-Up Option} --- add Rectified Adam as an
        alternative Phase~1 optimizer (\S\ref{sec:radam}).
\end{enumerate}

\medskip
\noindent
\textbf{Context.}  The original paper \cite{kiyani2025} tests five PDEs
(Burgers, Allen--Cahn, Kuramoto--Sivashinsky, Ginzburg--Landau, Stokes)
plus DeepONet.  \emph{Cahn--Hilliard is not among them}, making our
work a genuine extension of the optimizer family to fourth-order
conserved dynamics.  We ran 17~ablation experiments on the Brown
University OSCAR cluster (SLURM job~484517) varying optimizer variant,
warmup duration, Adam learning rate, batch size, power transform, and
Hessian scaling.


% ======================================================================
\section{Priority~1: Spectral Reference Solution}
\label{sec:spectral}
% ======================================================================

\subsection{Motivation}

The PINN training loss is \emph{not} a reliable proxy for solution
accuracy.  To quantify how well each trained PINN approximates the true
dynamics we need an independent, high-accuracy reference solution.
Because the Cahn--Hilliard equation is stiff (the biharmonic
$\nabla^4$ operator has eigenvalues $\sim k^4$), an explicit time
integrator would require impractically small time steps.  We therefore
use a fully spectral method.

\subsection{Method: ETDRK4}
\label{sec:etdrk4}

We implemented the \textbf{Exponential Time Differencing Runge--Kutta
scheme of order~4} (ETDRK4) following Kassam \& Trefethen
\cite{kassam2005}.  The idea is to treat the linear part of the PDE
exactly via matrix exponentials and integrate only the nonlinear part
numerically.

\paragraph{Semi-discrete form.}
In Fourier space the Cahn--Hilliard equation becomes
\begin{equation}
  \frac{d\hat{U}_{\bm{k}}}{dt}
  = \underbrace{-|\bm{k}|^2\bigl(|\bm{k}|^2 - a_2\bigr)}_{L_{\bm{k}}}\;
    \hat{U}_{\bm{k}}
  + \underbrace{\bigl(-|\bm{k}|^2 \cdot a_4\bigr)\,
    \widehat{U^3}_{\bm{k}}}_{N(\hat{U})_{\bm{k}}}
  \label{eq:fourier-ch}
\end{equation}
where $\bm{k} = (k_x, k_y)$ are the 2-D wavenumbers, $L_{\bm{k}}$ is
the (diagonal) linear operator, and $N$ contains the cubic nonlinearity
evaluated via FFT (dealiased by the $\tfrac{2}{3}$-rule implicitly
through the grid).

\paragraph{ETDRK4 coefficients.}
Following \cite{kassam2005}, we compute the coefficients $E$, $E_2$,
$a$, $b$, $c$ via contour integrals in the complex plane (32-point
circle), which avoids catastrophic cancellation that plagues direct
Taylor expansion for small $|L_{\bm{k}} \Delta t|$.

\paragraph{Parameters.}
\begin{itemize}[nosep]
  \item Grid: $128 \times 128$ (matching the PINN domain), periodic BCs
  \item Time step: $\Delta t = 0.005$, $t \in [0, 20]$, 101 snapshots
        saved (every 0.2 time units)
  \item Initial condition: identical to the PINN IC --- seed~42,
        $U_0 \sim \mathrm{Uniform}(-1,1)$ per grid point, Gaussian
        smoothed with $\sigma = 2.0$
\end{itemize}

\paragraph{Verification.}
We verified convergence by comparing $\Delta t = 0.005$ against
$\Delta t = 0.001$.  The maximum pointwise difference across all 101
snapshots is $\sim 10^{-13}$, i.e.\ machine precision for
\texttt{float64}.  Mass is perfectly conserved:
$\langle U \rangle(t{=}0) = \langle U \rangle(t{=}20) = -0.00685583$,
with zero drift.

\paragraph{Output.}
The solver writes \texttt{reference\_solution.npz} (12.1\,MB)
containing:
\begin{itemize}[nosep]
  \item \texttt{U}: shape $(101, 128, 128)$ --- solution snapshots
  \item \texttt{t}: shape $(101,)$ --- time values
  \item \texttt{x}, \texttt{y}: shape $(128,)$ --- spatial coordinates
\end{itemize}

\subsection{New file: \texttt{spectral\_solver.py}}

The solver is implemented in \texttt{spectral\_solver.py} ($\sim$240
lines).  Key functions:

\begin{itemize}[nosep]
  \item \texttt{generate\_ic()} --- reproduces the exact PINN initial
        condition (seed=42, $\sigma=2.0$, \texttt{Uniform}$[-1,1]$)
  \item \texttt{build\_wavenumbers()} --- constructs 2-D wavenumber
        grids with \texttt{fftfreq}
  \item \texttt{build\_etdrk4\_coefficients()} --- Kassam--Trefethen
        contour-integral coefficient computation
  \item \texttt{step\_etdrk4()} --- one ETDRK4 time step
  \item \texttt{solve\_cahn\_hilliard()} --- main driver with snapshot
        saving
\end{itemize}

\subsection{Relative $\ell_2$ Error Computation}

A second script, \texttt{compute\_l2\_errors.py} ($\sim$260 lines),
loads each of the 17~trained PINN models, evaluates them on the
reference grid at all 101~time snapshots, and computes the relative
$\ell_2$ error:
\begin{equation}
  \varepsilon_{\ell_2}(t_i)
  = \frac{\| U_{\text{PINN}}(t_i) - U_{\text{ref}}(t_i) \|_2}
         {\| U_{\text{ref}}(t_i) \|_2},
  \qquad
  \varepsilon_{\ell_2}^{\text{overall}}
  = \frac{\| U_{\text{PINN}} - U_{\text{ref}} \|_2}
         {\| U_{\text{ref}} \|_2}
  \label{eq:l2-error}
\end{equation}
where norms are taken over all grid points (and all times for the
overall error).  Results are saved to each experiment's
\texttt{summary.json} and aggregated in
\texttt{predictions/l2\_errors.json}.


\subsection{Results: $\ell_2$ Error Table}
\label{sec:l2-table}

Table~\ref{tab:l2-errors} shows the relative $\ell_2$ errors for all
17~experiments, sorted by overall error (best to worst).

\begin{table}[ht]
\centering
\caption{Relative $\ell_2$ errors against the spectral reference
  solution.  $\varepsilon_{\ell_2}$ is defined in
  Eq.~\eqref{eq:l2-error}.  ``Loss'' is the final PINN training loss.
  All errors are $\mathcal{O}(1)$, indicating that \emph{no experiment
  accurately captures the true Cahn--Hilliard dynamics}.}
\label{tab:l2-errors}
\small
\begin{tabular}{@{}llccccr@{}}
\toprule
\textbf{Experiment} & \textbf{Optimizer} &
  $\bm{\varepsilon_{\ell_2}^{\text{overall}}}$ &
  $\bm{\varepsilon_{\ell_2}(t{=}0)}$ &
  $\bm{\varepsilon_{\ell_2}(t{=}10)}$ &
  $\bm{\varepsilon_{\ell_2}(t{=}20)}$ &
  \textbf{Loss} \\
\midrule
bfgs\_batch\_200        & SSBroyden2  & \textbf{0.8545} & 0.5506 & 1.0265 & 0.9918 & $3.30 \times 10^{-4}$ \\
power\_4                & SSBroyden2  & 0.9451 & 0.7016 & 1.0004 & 1.0064 & $3.82 \times 10^{-3}$ \\
warmup\_500             & SSBroyden2  & 0.9487 & 0.6990 & 1.0752 & 1.0830 & $6.99 \times 10^{-4}$ \\
adam\_lr\_high           & SSBroyden2  & 0.9510 & 0.6931 & 1.1180 & 1.2044 & $2.36 \times 10^{-4}$ \\
bfgs\_vanilla           & BFGS\_scipy & 0.9748 & 0.6780 & 1.1756 & 1.1322 & $1.45 \times 10^{-3}$ \\
ssbroyden1              & SSBroyden1  & 0.9860 & 0.9953 & 0.9779 & 0.9710 & $1.63 \times 10^{0\phantom{-}}$ \\
adam\_only               & none        & 0.9997 & 0.9268 & 1.0522 & 1.0661 & $2.72 \times 10^{-2}$ \\
power\_2                & SSBroyden2  & 1.0613 & 0.8220 & 1.1541 & 1.2377 & $1.26 \times 10^{-2}$ \\
lbfgs                   & L-BFGS-B    & 1.0620 & 0.7814 & 1.2138 & 1.1164 & $1.05 \times 10^{-2}$ \\
ssbfgs\_ab              & SSBFGS\_AB  & 1.0712 & 0.6699 & 1.3526 & 1.3237 & $3.40 \times 10^{-4}$ \\
canonical               & SSBroyden2  & 1.0784 & 0.7397 & 1.1899 & 1.2020 & $5.10 \times 10^{-3}$ \\
bfgs\_batch\_1000       & SSBroyden2  & 1.1783 & 0.9575 & 1.3173 & 1.5134 & $2.65 \times 10^{-2}$ \\
warmup\_5000            & SSBroyden2  & 1.1942 & 0.6206 & 1.4896 & 1.9578 & $2.86 \times 10^{-4}$ \\
hessian\_scaled         & SSBroyden2  & 1.2956 & 0.7842 & 1.5198 & 1.7877 & $9.83 \times 10^{-3}$ \\
ssbfgs\_ol              & SSBFGS\_OL  & 1.3490 & 0.7922 & 1.6351 & 2.2455 & $1.44 \times 10^{-3}$ \\
warmup\_0               & SSBroyden2  & 1.6586 & 0.6337 & 2.1470 & 3.0792 & $3.89 \times 10^{-4}$ \\
adam\_lr\_low            & SSBroyden2  & 1.6804 & 0.6156 & 2.1796 & 3.1081 & $4.75 \times 10^{-4}$ \\
\bottomrule
\end{tabular}
\end{table}

\subsection{Key Findings from the $\ell_2$ Analysis}

\begin{enumerate}
  \item \textbf{All errors are $\mathcal{O}(1)$.}  The best overall
        $\ell_2$ error is 0.85 (\texttt{bfgs\_batch\_200}), meaning
        even the best model deviates from the reference by 85\% in a
        relative $\ell_2$ sense.  No PINN accurately solves the PDE.

  \item \textbf{Loss does not predict accuracy.}  The experiment with
        the \emph{lowest} training loss (\texttt{adam\_lr\_high},
        $2.36 \times 10^{-4}$) ranks only 4th in $\ell_2$ error
        (0.95).  Conversely, \texttt{ssbroyden1} has catastrophic loss
        ($1.63$) yet its $\ell_2 \approx 0.99$ is not the worst ---
        because its weights barely changed from initialization, so its
        output resembles the initial condition everywhere.

  \item \textbf{Errors grow with time.}  At $t{=}0$ all models have
        $\varepsilon_{\ell_2} \approx 0.6$--$1.0$, but by $t{=}20$
        the worst two experiments (\texttt{warmup\_0},
        \texttt{adam\_lr\_low}) reach $\varepsilon_{\ell_2} > 3.0$.
        This progressive degradation suggests the 1{,}361-parameter
        network cannot represent the complex phase-separation dynamics
        that develop at later times.

  \item \textbf{Implication.}  The [3,\,20,\,20,\,20,\,20,\,1] tanh
        architecture with 1{,}361 parameters is likely too small for
        the Cahn--Hilliard equation on a $128 \times 128$ domain over
        $t \in [0, 20]$.  This motivates Priority~3 (deeper
        architectures) and Priority~4 (time-marching / sequential
        PINNs).
\end{enumerate}


% ======================================================================
\section{Priority~2: SSBroyden1 Bug Fix \& Line Search Audit}
\label{sec:ssbroyden1}
% ======================================================================

All changes were made in \texttt{\_optimize.py}, the patched SciPy
optimizer file from \cite{kiyani2025}.  Every modification is bracketed
with \texttt{\# --- DRP MODIFICATION START/END ---} comments and includes
the original code for easy revert.

\subsection{Bug 1: Debug Print Statement}
\label{sec:bug-print}

The original SSBroyden1 block contained:
\begin{lstlisting}[caption={Original debug print (removed)}]
# ORIGINAL (line ~1693 in paper code):
print(ak, thetak, tauk, phik)
\end{lstlisting}
This executed \emph{every single iteration}, printing four floats per
line.  Over 10{,}000 BFGS iterations this floods stdout, slows SLURM
log writes, and can cause I/O bottlenecks on the cluster.

\textbf{Fix:} Removed.

\subsection{Bug 2: Missing NaN/Inf Guards in SSBroyden1}
\label{sec:bug-nan}

The SSBroyden1 Hessian update involves a chain of intermediate
quantities that can produce division-by-zero, NaN, or Inf:

\begin{center}
\renewcommand{\arraystretch}{1.3}
\begin{tabular}{@{}lll@{}}
\toprule
\textbf{Quantity} & \textbf{Formula} & \textbf{Singularity} \\
\midrule
$\theta_k^{\mathrm{SR1}}$ & $1 / (1 - b_k)$ & $b_k \to 1$ \\
$\theta_k^{-}$ & $(\rho_k^{-} - 1) / a_k$ & $a_k \to 0$ \\
$\theta_k^{+}$ & $1 / \rho_k^{-}$ & $\rho_k^{-} \to 0$ \\
$\rho_k^{-}$ & $\min\!\bigl(1,\; h_k(1 - \sqrt{|a_k|/(1+a_k)})\bigr)$ & $1 + a_k \le 0$ \\
$v_k$ & $s_k \rho_k - H_k y_k / (y_k^T H_k y_k)$ & $y_k^T H_k y_k \to 0$ \\
\bottomrule
\end{tabular}
\end{center}

Once \emph{any} of these quantities becomes NaN or Inf, it corrupts
the Hessian approximation $H_k$, and all subsequent iterations diverge.
This explains the catastrophic failure of \texttt{ssbroyden1} (final
loss = 1.63, essentially untrained).

\textbf{Fix:} We introduced an epsilon guard \texttt{\_EPS\_DIV =
1e-30} and a boolean \texttt{\_fallback} flag.  Each intermediate
computation checks for near-zero denominators and non-finite results.
If any check fails, \texttt{\_fallback = True} and the iteration uses
a \textbf{standard BFGS update} instead of the SSBroyden1 formula.
This ensures the Hessian remains well-conditioned even when the
SSBroyden1-specific intermediates degenerate.

The complete rewritten block is at lines~1618--1753 of
\texttt{\_optimize.py}, with the original code preserved verbatim in
comments.

\subsection{Bug 3: Negative Curvature Guard (All Variants)}
\label{sec:bug-curvature}

The curvature condition for quasi-Newton methods requires $y_k^T s_k >
0$, where $y_k = g_{k+1} - g_k$ and $s_k = x_{k+1} - x_k$.  The
quantity \texttt{rhok\_inv} $= y_k^T s_k$ is used to compute $\rho_k =
1 / (y_k^T s_k)$.

\textbf{Original code (all variants):}
\begin{lstlisting}[caption={Original rhok\_inv guard}]
rhok_inv = np.dot(yk, sk)
if rhok_inv == 0.:
    rhok = 1000.0
else:
    rhok = 1. / rhok_inv
\end{lstlisting}

This only guards against exact zero, but not \textbf{negative} values.
A negative $y_k^T s_k$ (violated curvature condition) can occur when the
Strong Wolfe line search accepts a step that does not perfectly satisfy
the conditions due to numerical noise.  A negative $\rho_k$ would
produce a descent direction that \emph{ascends}, causing divergence.

\textbf{Fix:}
\begin{lstlisting}[caption={New negative curvature guard (lines 1555--1568)}]
rhok_inv = np.dot(yk, sk)
if rhok_inv == 0.:
    rhok = 1000.0
    if disp:
        msg = "Divide-by-zero encountered: rhok assumed large"
        _print_success_message_or_warn(True, msg)
# --- DRP MODIFICATION START (negative curvature guard) ---
elif rhok_inv < 0.:
    rhok = 1000.0
    if disp:
        msg = "Negative curvature detected: rhok assumed large"
        _print_success_message_or_warn(True, msg)
# --- DRP MODIFICATION END ---
else:
    rhok = 1. / rhok_inv
\end{lstlisting}

Setting $\rho_k = 1000$ when $y_k^T s_k < 0$ effectively makes the
BFGS update vanishingly small (since the update terms scale with
$\rho_k$ and $1000 \gg 1/\|y\|\|s\|$ in practice), keeping $H_k$
nearly unchanged --- equivalent to a steepest-descent step.  This
guard affects \textbf{all six quasi-Newton variants} (BFGS,
BFGS\_scipy, SSBFGS\_OL, SSBFGS\_AB, SSBroyden1, SSBroyden2) since
the $\rho_k$ computation precedes the variant-specific blocks.

\subsection{Line Search Audit}
\label{sec:line-search}

We verified that the line search parameters match the paper exactly:
\begin{itemize}[nosep]
  \item \textbf{Strong Wolfe conditions} with $c_1 = 10^{-4}$,
        $c_2 = 0.9$ (function signature, line~1415)
  \item $c_1$ and $c_2$ are correctly passed to
        \texttt{\_line\_search\_wolfe12} (line~1505)
  \item No changes were needed to the line search itself
\end{itemize}


\subsection{Summary of \texttt{\_optimize.py} Changes}

\begin{table}[ht]
\centering
\caption{All modifications to \texttt{\_optimize.py}.  Each is marked
  with \texttt{DRP MODIFICATION} comments in the source.}
\label{tab:optimize-changes}
\begin{tabular}{@{}clp{7.5cm}@{}}
\toprule
\textbf{\#} & \textbf{Lines} & \textbf{Description} \\
\midrule
1 & 1555--1568 & Added \texttt{elif rhok\_inv < 0} guard (negative
    curvature).  Affects all 6 quasi-Newton variants. \\
2 & 1618--1753 & Complete rewrite of SSBroyden1 block: removed debug
    print, added epsilon guards on 5~singularities, added BFGS
    fallback.  Original code preserved in comments. \\
\bottomrule
\end{tabular}
\end{table}


% ======================================================================
\section{Priority~5: RAdam Warm-Up}
\label{sec:radam}
% ======================================================================

\subsection{Motivation}

The Phase~1 warm-up uses Adam \cite{kingma2015} to pre-train the
network before handing off to a quasi-Newton optimizer.  Standard Adam
can be sensitive to the learning rate in the first few hundred steps
because the exponential moving average of squared gradients (the
variance estimate) has high bias when few samples have been seen.

\textbf{Rectified Adam} (RAdam) \cite{liu2020} addresses this by
dynamically computing the variance of the adaptive learning rate and
turning off momentum correction when the variance is too high.  This
acts as a \emph{built-in warm-up} mechanism, eliminating the need for
a manually tuned learning rate schedule in early training.

\subsection{Changes to \texttt{run\_experiment.py}}

Two modifications were made:

\paragraph{1. Config parsing (line~97).}
\begin{lstlisting}[caption={New config key for optimizer selection}]
adam_optimizer = adam_cfg.get("optimizer", "adam").lower()
\end{lstlisting}
The key \texttt{optimizer} is optional and defaults to \texttt{"adam"},
so all 17~existing config files are backward-compatible without any
changes.

\paragraph{2. Optimizer dispatch (lines~453--458).}
\begin{lstlisting}[caption={RAdam / Adam dispatch}]
if Nepochs_ADAM > 0:
    _optim_cls = {"adam": torch.optim.Adam,
                  "radam": torch.optim.RAdam}
    if adam_optimizer not in _optim_cls:
        raise ValueError(
            f"Unknown optimizer '{adam_optimizer}'; "
            f"use 'adam' or 'radam'"
        )
    optimizer = _optim_cls[adam_optimizer](
        net.parameters(), lr=adam_lr,
        betas=adam_betas, eps=adam_eps
    )
\end{lstlisting}

\subsection{New config: \texttt{configs/radam\_warmup.yaml}}

A new YAML config was created, identical to the canonical SSBroyden2
configuration except for one line:
\begin{lstlisting}[language={},caption={RAdam config change}]
training:
  adam:
    optimizer: "radam"   # <-- only change
    epochs: 2000
    lr: 0.001
    ...
\end{lstlisting}

All other parameters (network architecture, sampling, loss weights,
BFGS variant, etc.) are identical to the canonical config, isolating
the effect of RAdam vs.\ Adam.

\subsection{Verification}

\begin{itemize}[nosep]
  \item \texttt{torch.optim.RAdam} confirmed available in PyTorch~2.6.0
  \item Canonical config correctly defaults to Adam (no
        \texttt{optimizer} key $\Rightarrow$ \texttt{"adam"})
  \item RAdam config correctly selects RAdam
  \item Syntax check passed (no import or parse errors)
\end{itemize}


% ======================================================================
\section{Files Created or Modified}
\label{sec:files}
% ======================================================================

\begin{table}[ht]
\centering
\caption{Complete inventory of file changes.}
\label{tab:files}
\begin{tabular}{@{}llp{7cm}@{}}
\toprule
\textbf{File} & \textbf{Status} & \textbf{Description} \\
\midrule
\texttt{spectral\_solver.py}        & \textcolor{diffadd}{NEW} &
  ETDRK4 spectral solver for Cahn--Hilliard ($\sim$240 lines) \\
\texttt{compute\_l2\_errors.py}     & \textcolor{diffadd}{NEW} &
  Evaluates 17 PINNs against reference, computes $\ell_2$ errors
  ($\sim$260 lines) \\
\texttt{reference\_solution.npz}    & \textcolor{diffadd}{NEW} &
  Spectral reference (12.1\,MB, 101 snapshots, $128 \times 128$) \\
\texttt{predictions/l2\_errors.json}& \textcolor{diffadd}{NEW} &
  Aggregated $\ell_2$ errors for all 17 experiments \\
\texttt{configs/radam\_warmup.yaml} & \textcolor{diffadd}{NEW} &
  RAdam warm-up ablation config \\
\texttt{\_optimize.py}              & \textcolor{orange}{MOD} &
  2 changes: negative curvature guard (lines 1555--1568), SSBroyden1
  rewrite (lines 1618--1753) \\
\texttt{run\_experiment.py}         & \textcolor{orange}{MOD} &
  2 changes: RAdam config parsing (line 97), optimizer dispatch
  (lines 453--458) \\
\bottomrule
\end{tabular}
\end{table}


% ======================================================================
\section{Next Steps}
\label{sec:next}
% ======================================================================

The following priorities remain for the next phase of work:

\begin{enumerate}
  \item \textbf{Re-run on Oscar.}  Push the fixed
        \texttt{\_optimize.py}, updated \texttt{run\_experiment.py},
        and new \texttt{radam\_warmup.yaml} to Oscar.  Re-run
        \texttt{ssbroyden1} to verify the fix and run the RAdam
        experiment.

  \item \textbf{Priority~3: Deeper Architectures.}  The
        $\mathcal{O}(1)$ $\ell_2$ errors strongly suggest the
        $[3,20,20,20,20,1]$ network (1{,}361 parameters) is
        insufficient for the Cahn--Hilliard equation.  Test wider
        (e.g.\ $[3,64,64,64,64,1]$) and deeper (e.g.\ 6--8~hidden
        layers) architectures.

  \item \textbf{Priority~4: Time-Marching.}  Train sequential PINNs
        on short time windows (e.g.\ $[0,2], [2,4], \ldots, [18,20]$)
        and stitch the solutions together.  This avoids the need for a
        single network to represent 20 time units of complex dynamics.

  \item \textbf{Updated report and plots.}  Once new experiments
        complete, update the analysis plots and the experiment report
        with $\ell_2$ error comparisons.
\end{enumerate}


% ======================================================================
% References
% ======================================================================
\bibliographystyle{plain}
\begin{thebibliography}{9}

\bibitem{kiyani2025}
  R.~Kiyani, O.~Yasa, M.~Fluid, et al.,
  ``Optimizing the Optimizer for Data-Driven PINNs and KANs,''
  \emph{Comput.\ Methods Appl.\ Mech.\ Engrg.}, vol.~446, 118308, 2025.
  \href{https://arxiv.org/abs/2501.16371}{arXiv:2501.16371}.

\bibitem{kassam2005}
  A.-K.~Kassam and L.~N.~Trefethen,
  ``Fourth-order time-stepping for stiff PDEs,''
  \emph{SIAM J.\ Sci.\ Comput.}, vol.~26, no.~4, pp.~1214--1233, 2005.

\bibitem{kingma2015}
  D.~P.~Kingma and J.~Ba,
  ``Adam: A Method for Stochastic Optimization,''
  \emph{Proc.\ ICLR}, 2015.
  \href{https://arxiv.org/abs/1412.6980}{arXiv:1412.6980}.

\bibitem{liu2020}
  L.~Liu, H.~Jiang, P.~He, et al.,
  ``On the Variance of the Adaptive Learning Rate and Beyond,''
  \emph{Proc.\ ICLR}, 2020.
  \href{https://arxiv.org/abs/1908.03265}{arXiv:1908.03265}.

\end{thebibliography}

\end{document}
